% ============================================================
% common/macros.tex —— 自定义命令与环境（题号、解析头、提示框等）
% 与合并书 math-gaokao-2017-2026/preamble.tex 的自定义段保持一致，
% 保证内容片段（questions.tex / answers.tex）在两个项目间可互换。
% ============================================================

% ---------- 试卷标识（壳文件 \renewcommand 设置，用于 label 命名空间）----------
\newcommand{\gkpaperid}{none}

% ---------- 题号 ----------
% 试题部分的题号：加粗编号 + 自动 \label{q:<卷id>:<题号>}
\newcommand{\gkti}[1]{\par\medskip\noindent
  \textbf{#1．}\phantomsection\label{q:\gkpaperid:#1}\ignorespaces}

% 解析部分的题头：小标题 + 回链到试题页
% （仅在题目与解析同文档的版本中回链有效；试题版不含解析头，无悬空引用）
\newcommand{\gksolhead}[2]{%
  \Needspace{3\baselineskip}\subsubsection*{第 #1 题\if\relax\detokenize{#2}\relax\else（#2）\fi}%
  \phantomsection\label{s:\gkpaperid:#1}%
  {\small\color{gray}（题目见第~\pageref{q:\gkpaperid:#1}~页）}\par\nopagebreak}

% ---------- 填空题空位 ----------
\newcommand{\gkblank}{\underline{\hspace{2.2em}}}

% ---------- 提示框 ----------
% 卷种提示 / 解析归属声明（醒目，浅黄）
\newtcolorbox{gknotice}{colback=yellow!8, colframe=orange!60!black,
  boxrule=0.5pt, arc=1.5pt, left=6pt, right=6pt, top=4pt, bottom=4pt,
  before skip=8pt, after skip=8pt}

% 方法提示 / 注（浅蓝，用于解析后的点评与题量说明）
\newtcolorbox{gktip}{colback=blue!4, colframe=blue!35!black,
  boxrule=0.4pt, arc=1.5pt, left=6pt, right=6pt, top=4pt, bottom=4pt,
  before skip=6pt, after skip=6pt, fontupper=\small}

% ---------- 答案行 ----------
% 解析中独立成段的"答案：X"用带底色的行呈现
\newcommand{\gkanswer}[1]{\par\smallskip\noindent
  \tcbox[width=\linewidth, tcbox width=auto limited,
         colback=green!6, colframe=green!35!black, boxrule=0.4pt,
         arc=1.5pt, left=4pt, right=4pt, top=2pt, bottom=2pt,
         nobeforeafter]{\textbf{#1}}\par\smallskip}

% ---------- 结构标题 ----------
% 结构标题后通常紧跟一个带 0.5\textheight 起始预留的解析题块；0.4 只够放标题，
% 会把下一题整体推到下一页，留下孤立的「二、填空题」等页尾标题。
\newcommand{\gkpart}[1]{\Needspace{0.6\textheight}\subsection*{#1}\nopagebreak}     % 标题 + 首题
% （一）必考题 / 解析小问标题：块式加粗段，可自动换行（\paragraph 的 run-in 版式会溢出）
\newcommand{\gksubpart}[1]{\Needspace{6\baselineskip}\par\medskip\noindent\textbf{#1}\par\nopagebreak\smallskip}

% 罗马数字Ⅰ-Ⅻ（U+2160）与圈号①②③（U+2460）划入 CJK 字符类，
% 用中文字体渲染——西文字体 Latin Modern 缺这些字形，否则整段静默丢字
\xeCJKDeclareCharClass{CJK}{"2160 -> "217F, "2460 -> "24FF}
